\documentclass[11pt]{article}
\usepackage[paperwidth=11in,paperheight=8.5in,margin=0.25in]{geometry}
\usepackage[T1]{fontenc}
\usepackage[utf8]{inputenc}
\usepackage[english]{babel}
\usepackage{lmodern}
\usepackage{microtype}
\usepackage{amsmath,amssymb}
\usepackage{xcolor}
\usepackage{tikz}
\usetikzlibrary{calc}

\hyphenpenalty=8000
\exhyphenpenalty=8000
\emergencystretch=2em

\definecolor{QCBlue}{HTML}{1F4E79}
\definecolor{QCGreen}{HTML}{2E6B5F}
\definecolor{QCOrange}{HTML}{B65B2A}
\definecolor{QCPurple}{HTML}{4B3F72}
\definecolor{QCLight}{HTML}{F4F6F8}
\definecolor{QCBorder}{HTML}{D0D6DC}
\definecolor{QCText}{HTML}{111111}

\newlength{\Margin} \setlength{\Margin}{0.25in}
\newlength{\Gap} \setlength{\Gap}{0.1in}
\newlength{\TitleH} \setlength{\TitleH}{0.8in}
\newlength{\HeaderH} \setlength{\HeaderH}{0.38in}
\newlength{\Pad} \setlength{\Pad}{0.1in}
\newlength{\UsableW}
\newlength{\UsableH}
\newlength{\CellW}
\newlength{\CellH}
\setlength{\UsableW}{\dimexpr\paperwidth-2\Margin\relax}
\setlength{\UsableH}{\dimexpr\paperheight-2\Margin\relax}
\setlength{\CellW}{5.2in}
\setlength{\CellH}{3.5in}

\newcommand{\QuadBox}[4]{%
  \node[anchor=north west, draw=QCBorder, line width=0.6pt, rounded corners=10pt,
        fill=white, minimum width=\CellW, minimum height=\CellH, inner sep=0pt] (B) at #1 {};
  \fill[#2, rounded corners=10pt]
    ([xshift=0pt,yshift=0pt]B.north west) rectangle
    ([xshift=0pt,yshift=-\HeaderH]B.north east);
  \node[anchor=west, text=white, font=\bfseries] at ([xshift=0.18in,yshift=-0.22in]B.north west) {#3};
  \node[anchor=north west, text=QCText] at ([xshift=\Pad,yshift=-\HeaderH-\Pad]B.north west) {%
    \begin{minipage}[t]{\dimexpr\CellW-2\Pad\relax}
      \raggedright\small #4
    \end{minipage}
  };
}

\pagestyle{empty}
\begin{document}
\begin{tikzpicture}[remember picture,overlay]
\coordinate (NW) at ([xshift=\Margin,yshift=-\Margin]current page.north west);
\node[anchor=north west, draw=QCBorder, line width=0.6pt, rounded corners=12pt,
      fill=QCLight, minimum width=\UsableW, minimum height=\TitleH, inner sep=0pt] (T) at (NW) {};
\node[anchor=north west] at ([xshift=0.22in,yshift=-0.22in]T.north west) {%
  \begin{minipage}[t]{\dimexpr\UsableW-0.44in\relax}
    {\Huge\bfseries \textcolor{QCBlue}{Poisson's Equation Formula Sheet}}\par\vspace{2pt}
    {\normalsize \textbf{Diego Juarez}\hfill \textbf{Computational Electrodynamics}\hfill \textbf{\today}}
  \end{minipage}
};
\coordinate (GridNW) at ([yshift=-0.15in]T.south west);
\coordinate (TL) at (GridNW);
\coordinate (TR) at ([xshift=\CellW+\Gap]GridNW);
\coordinate (BL) at ([yshift=-\CellH-\Gap]GridNW);
\coordinate (BR) at ([xshift=\CellW+\Gap,yshift=-\CellH-\Gap]GridNW);

\QuadBox{(TL)}{QCBlue}{1) Governing Equations}{
\textbf{Electrostatics}
\[
\mathbf{E}=-\nabla\Phi, \qquad \nabla\cdot\mathbf{E}=\frac{\rho}{\varepsilon_0}
\]
\[
\boxed{\nabla^2\Phi=-\frac{\rho}{\varepsilon_0}}
\]
Charge-free region:
\[
\rho=0 \Longrightarrow \nabla^2\Phi=0
\]
Linearity:
\[
\Phi=a\Phi_1+b\Phi_2 \Longrightarrow \rho=a\rho_1+b\rho_2
\]
Cartesian Laplacian:
\[
\nabla^2=\frac{\partial^2}{\partial x^2}+\frac{\partial^2}{\partial y^2}+\frac{\partial^2}{\partial z^2}
\]
Spherical symmetry:
\[
\nabla^2\Phi=\frac{1}{r^2}\frac{d}{dr}\left(r^2\frac{d\Phi}{dr}\right)
\]
1D form:
\[
\frac{d^2\Phi}{dx^2}=-\frac{\rho(x)}{\varepsilon_0}
\]
}

\QuadBox{(TR)}{QCGreen}{2) Integral and Symmetry Results}{
Free-space solution:
\[
\boxed{\Phi(\mathbf r)=\frac{1}{4\pi\varepsilon_0}\int \frac{\rho(\mathbf r')}{|\mathbf r-\mathbf r'|}\,d^3r'}
\]
Point charge:
\[
\rho(\mathbf r)=q\delta^{(3)}(\mathbf r), \qquad
\Phi=\frac{q}{4\pi\varepsilon_0 r}
\]
Radial field from Gauss's law:
\[
E_r(r)\,4\pi r^2=\frac{Q_{\text{enc}}(r)}{\varepsilon_0}
\]
Potential from field:
\[
\Phi(r_2)-\Phi(r_1)=-\int_{r_1}^{r_2} E_r(r)\,dr
\]
Uniform sphere $(r<R)$:
\[
E_r=\frac{\rho_0 r}{3\varepsilon_0}
\]
Uniform sphere $(r>R)$:
\[
E_r=\frac{1}{4\pi\varepsilon_0}\frac{Q}{r^2}, \qquad Q=\frac{4\pi R^3}{3}\rho_0
\]
}

\QuadBox{(BL)}{QCOrange}{3) Boundary and Uniqueness Facts}{
\textbf{Dirichlet data:} specify $\Phi$ on the boundary.

\textbf{Neumann data:} specify $\partial \Phi/\partial n$ on the boundary.

At an interface with no surface charge:
\[
\Phi_1=\Phi_2, \qquad
\left.\frac{\partial \Phi_1}{\partial n}\right|_S=
\left.\frac{\partial \Phi_2}{\partial n}\right|_S
\]
With surface charge density $\sigma$:
\[
E_{2n}-E_{1n}=\frac{\sigma}{\varepsilon_0}
\]
Uniqueness idea:
\[
\nabla^2(\Phi_1-\Phi_2)=0
\]
and zero boundary data imply $\Phi_1=\Phi_2$.

Physical reminder:

Potential in a region can be nontrivial even when $\rho=0$ there, because boundary values encode external sources.
}

\QuadBox{(BR)}{QCPurple}{4) Useful Solvers and Checks}{
\textbf{Finite difference}
\[
\frac{\Phi_{i+1}-2\Phi_i+\Phi_{i-1}}{h^2}\approx -\frac{\rho_i}{\varepsilon_0}
\]
Second-order truncation error:
\[
\mathcal O(h^2)
\]
\textbf{1D jump from delta source}
\[
\Phi''(x)=-\frac{q}{\varepsilon_0}\delta(x-a)
\]
\[
\Phi(a^+)=\Phi(a^-), \qquad
\Phi'(a^+)-\Phi'(a^-)=-\frac{q}{\varepsilon_0}
\]
\textbf{Sign check}

Positive $\rho$ gives negative curvature of $\Phi$ in 1D.

\textbf{Units}
\[
[\Phi]=\mathrm V,\quad [\rho]=\mathrm{C\,m^{-3}},\quad [\varepsilon_0]=\mathrm{C\,V^{-1}m^{-1}}
\]
\[
[\nabla^2\Phi]=\mathrm{V\,m^{-2}}=[\rho/\varepsilon_0]
\]
}
\end{tikzpicture}
\end{document}
